%%%%%
%
\section{Point de ros�e sachant l'humidit� relative}

\itemstyle{Op\'erateur~:} \operfmt{Td\_Hu}
\index{Td\_Hu}

%%%%%OBJET%%%%
\itemstyle{Objet~:}
 Calcul de la temp�rature du point de ros�e � partir de l'humidit�
relative � l'eau liquide et la temp�rature de l'air.

%%%%%ARGUMENTS%%%%
\bigskip
%%%%%ENTREES%%%%
\begin{tabarg}
\multicolumn{3}{l}{\titretab {Entr\'ees}}\\
\hline
$Hu$ & Vecteur, \texttt{float} & Humidit� relative (\%)\\
$Ts$ & Vecteur, \texttt{float} & Temp�rature de l'air (K ou Celsius)\\
\end{tabarg}

\medskip
%%%%%SORTIES%%%%
\begin{tabarg}
\multicolumn{3}{l}{\titretab {Sorties}}\\
\hline
$Td$ & Vecteur, \texttt{float} & Temp�rature du point de ros�e (Celsius)\\
\end{tabarg}


%%%%%FORMULE%%%%
\itemstyle{Formule}

%%Td_Hu
$Ts$ est converti en Kelvin, si besoin est, puis :
\[
e = \frac{Hu}{100} . e_w(Ts) \qquad \text{(d�finition de Hu)}
\]
et
\[
Td_K = e_{w}^{-1}(e)
\]
enfin, $Td = Td_K -273,15 $.


%%%%%AUTEUR%%%%
\itemstyle{Auteur~:}
CNRM/GMEI/TRAMM

%%%%%REFERENCE%%%%
\itemstyle{R�f�rences~:}

Les formules directes et r�ciproques de $ew(T)$ sont
celles de l'article Bolton-1980 \cite{Bolton-80} :
eq. 10, page 1047, d�crite en annexe 
\ref{annexe-ewbolton}.

L'OMM sp�cifie explicitement
que, aux temp�ratures inf�rieures � O�C,
l'humidit� relative doit
�tre calcul�e relativement � l'eau liquide.
(par exemple dans le
\emph{Guide des instruments et des m�thodes
d'observation m�t�orologique} \cite{OMM-N8-1996}, en
annexe 4.A et au paragraphe 12.5.1).
C'est pourquoi, ici, seule la formule
relative � l'eau liquide est employ�e pour
le calcul de $e$ � partir de Hu.
